%! Principal Component Analysis — Unit 7, Lesson 7.3 — Exit Ticket
\documentclass[10pt]{article}
\usepackage{linalg-article}
\usepackage{linalg-boxes}
\newcommand{\TallMath}[1]{$\displaystyle #1\rule[-1.4em]{0pt}{3.2em}$}
\newcommand{\uu}{\mathbf{u}}
\newcommand{\vv}{\mathbf{v}}
\newcommand{\zero}{\mathbf{0}}
\newcommand{\T}{^{\top}}

\begin{document}

\pageheader{Unit 7, Lesson 7.3}{Exit Ticket}
\namedateperiod

\vspace{0.15in}

No notes. Show your reasoning. A data set of four points has already been \textbf{centered}:
$(2,3),(3,2),(-2,-3),(-3,-2)$. Stacked as the rows of $A$, they give
$A\T A=\begin{bmatrix} 26 & 24\\ 24 & 26 \end{bmatrix}$.

\begin{enumerate}[leftmargin=1.8em, itemsep=15pt]
  \item \textbf{Find the principal components.} Solve $\det(A\T A-\lambda I)=0$ for $\lambda_1,\lambda_2$, then find PC1 (the unit eigenvector for the larger $\lambda$).
        \vspace{1.9cm}
  \item \textbf{How much spread on PC1?} The spread along a component is $\sigma^2=\lambda$. What fraction of the total spread does PC1 capture? Is summarizing each point by its PC1 coordinate a good idea here?
        \vspace{1.9cm}
  \item \textbf{Explain.} (a) Why do we \emph{center} the data before finding principal components? (b) What does keeping only PC1 let us do with the data? \writelines{2}
\end{enumerate}

\end{document}
